% 参数计数已在正文由 $2N_g-1$ 个可消去重相位一行完成；不另设推导框。

\begin{derivation}[从\eqnrefs{eq:ckm-rephasing-dp68}到\eqnrefs{eq:jarlskog-dp11}]
质量本征场仍允许保持质量项不变的矢量式重相位

\begin{equation*}
 u_i\to \ee^{\ii\alpha_i}u_i,
 \qquad
 d_j\to \ee^{\ii\beta_j}d_j.
\end{equation*}
带电流中的 CKM 元素因此变换为
\begin{equation*}
 V_{ij}\to
 \ee^{-\ii\alpha_i}V_{ij}\ee^{\ii\beta_j}.
\end{equation*}
考虑四元积
\begin{equation*}
 X_{ik;jl}\equiv V_{ij}V_{kl}V_{il}^*V_{kj}^*,
 \qquad i\neq k,
 \quad j\neq l.
\end{equation*}
四个因子分别携带相位
\begin{align*}
 V_{ij}:&\quad \ee^{-\ii\alpha_i}\ee^{+\ii\beta_j},\\
 V_{kl}:&\quad \ee^{-\ii\alpha_k}\ee^{+\ii\beta_l},\\
 V_{il}^*:&\quad \ee^{+\ii\alpha_i}\ee^{-\ii\beta_l},\\
 V_{kj}^*:&\quad \ee^{+\ii\alpha_k}\ee^{-\ii\beta_j}.
\end{align*}
相乘后 $\alpha_i,\alpha_k,\beta_j,\beta_l$ 各出现一次正号和一次负号，总相位严格等于 $1$。因此 $X_{ik;jl}$ 以及它的虚部都不依赖质量本征场的重相位选择；单个 CKM 元素的相位则没有这种性质。

取标准三角参数化 $V=R_{23}U_{13}(\delta)R_{12}$。计算\eqnrefs{eq:jarlskog-dp11} 可选择
\begin{equation*}
 J_{\rm CKM}
 =\Im\left(V_{ud}V_{cs}V_{us}^*V_{cd}^*\right).
\end{equation*}
所需四个矩阵元为
\begin{align*}
 V_{ud}&=c_{12}c_{13},\\
 V_{us}&=s_{12}c_{13},\\
 V_{cd}&=-s_{12}c_{23}-c_{12}s_{23}s_{13}\ee^{\ii\delta},\\
 V_{cs}&=c_{12}c_{23}-s_{12}s_{23}s_{13}\ee^{\ii\delta}.
\end{align*}
前两个矩阵元为实数，因此
\begin{align*}
 J_{\rm CKM}
 =c_{12}s_{12}c_{13}^2\,
 \Im\Big[&
 (c_{12}c_{23}-s_{12}s_{23}s_{13}\ee^{\ii\delta})\\
 &\times(-s_{12}c_{23}-c_{12}s_{23}s_{13}\ee^{-\ii\delta})
 \Big].
\end{align*}
把括号展开：
\begin{align*}
 {}&-c_{12}s_{12}c_{23}^2
 -c_{12}^2c_{23}s_{23}s_{13}\ee^{-\ii\delta}\\
 &\qquad
 +s_{12}^2c_{23}s_{23}s_{13}\ee^{\ii\delta}
 +c_{12}s_{12}s_{23}^2s_{13}^2.
\end{align*}
第一项和最后一项为实数；中间两项的虚部为
\begin{align*}
 &c_{12}^2c_{23}s_{23}s_{13}\sin\delta
 +s_{12}^2c_{23}s_{23}s_{13}\sin\delta\\
 &=(c_{12}^2+s_{12}^2)c_{23}s_{23}s_{13}\sin\delta\\
 &=c_{23}s_{23}s_{13}\sin\delta.
\end{align*}
所以
\begin{equation*}
 J_{\rm CKM}
 =c_{12}c_{23}c_{13}^2s_{12}s_{23}s_{13}\sin\delta,
\end{equation*}
即\eqnrefs{eq:jarlskog-dp11}。

三代幺正性进一步固定所有异行异列四元积的相对符号。定义
\begin{equation*}
 \mathcal J_{ij;\alpha\beta}
 \equiv\Im\!\left(
 V_{i\alpha}V_{j\beta}V_{i\beta}^*V_{j\alpha}^*
 \right),
 \qquad i\neq j,\quad \alpha\neq\beta.
\end{equation*}
交换两行或两列都会把括号内的复数变为其复共轭，因此
\begin{equation*}
 \mathcal J_{ji;\alpha\beta}
 =-\mathcal J_{ij;\alpha\beta},
 \qquad
 \mathcal J_{ij;\beta\alpha}
 =-\mathcal J_{ij;\alpha\beta}.
\end{equation*}
再固定一对不同的行 $i,j$ 与一列 $\alpha$。行正交性给
\begin{equation*}
 \sum_{\lambda=1}^{3}V_{i\lambda}V_{j\lambda}^*=0.
\end{equation*}
乘以 $V_{i\alpha}^*V_{j\alpha}$ 并取虚部，$\lambda=\alpha$ 项为实数，于是其余两列 $\beta,\gamma$ 满足
\begin{equation*}
 \mathcal J_{ij;\beta\alpha}
 +\mathcal J_{ij;\gamma\alpha}=0.
\end{equation*}
结合列交换的反对称性，固定行对时三个不同列对的四元积虚部具有相同绝对值，并按列指标的奇偶排列改变符号。

为了比较不同的行对，固定两列 $\alpha\neq\beta$ 与一行 $i$。列正交性给
\begin{equation*}
 \sum_{r=1}^{3}V_{r\alpha}V_{r\beta}^*=0.
\end{equation*}
乘以 $V_{i\alpha}^*V_{i\beta}$ 并取虚部，$r=i$ 项为
$|V_{i\alpha}|^2|V_{i\beta}|^2$，其虚部为零，因此剩余两行 $j,k$ 满足
\begin{equation*}
 \Im\!\left(
 V_{j\alpha}V_{j\beta}^*
 V_{i\alpha}^*V_{i\beta}
 \right)
 +
 \Im\!\left(
 V_{k\alpha}V_{k\beta}^*
 V_{i\alpha}^*V_{i\beta}
 \right)=0.
\end{equation*}
把每一项重新排列成 $\mathcal J$ 的定义，并使用行交换的反对称性，就得到固定列对时不同的行对具有相同绝对值，并按行指标排列改变符号。以
\begin{equation*}
 \mathcal J_{12;12}
 =\Im(V_{ud}V_{cs}V_{us}^*V_{cd}^*)
 \equiv J_{\rm CKM}
\end{equation*}
固定总体符号，便得到三代幺正矩阵的恒等式
\begin{equation*}
 \mathcal J_{ij;\alpha\beta}
 =J_{\rm CKM}\,\epsilon_{ijk}\epsilon_{\alpha\beta\gamma},
\end{equation*}
其中 $k$ 与 $\gamma$ 分别是剩余的第三个行、列指标。这就是三代幺正矩阵的一般四元积恒等式。

\end{derivation}

\begin{derivation}[从\eqnrefs{eq:jarlskog-dp11}到\eqnrefs{eq:jarlskog-determinant-dp11}]
定义两个厄米矩阵

\begin{equation*}
 H_u\equiv Y_uY_u^\dagger,
 \qquad
 H_d\equiv Y_dY_d^\dagger.
\end{equation*}
在共同左手弱基变换 $Q_L\to U_QQ_L$ 下，
\begin{equation*}
 H_u\to U_QH_uU_Q^\dagger,
 \qquad
 H_d\to U_QH_dU_Q^\dagger.
\end{equation*}
因此对易子
\begin{equation*}
 C\equiv[H_u,H_d]
\end{equation*}
只作相似变换 $C\to U_QCU_Q^\dagger$，所以 $\det C$ 与弱基选择无关。又因为 $H_u,H_d$ 厄米，
\begin{equation*}
 C^\dagger=-C,
\end{equation*}
三代时 $\det C$ 必为纯虚数。

选择上型夸克质量基，并简记
\begin{equation*}
 a_i\equiv y_{u_i}^2,
 \qquad
 b_\alpha\equiv y_{d_\alpha}^2.
\end{equation*}
于是
\begin{equation*}
 H_u=\operatorname{diag}(a_1,a_2,a_3),
 \qquad
 H_d=V\operatorname{diag}(b_1,b_2,b_3)V^\dagger,
\end{equation*}
其中 $V=V_{\rm CKM}$。对易子的矩阵元为
\begin{align*}
 C_{ij}
 &=(a_i-a_j)(H_d)_{ij},\\
 (H_d)_{ij}
 &=\sum_{\alpha=1}^{3}V_{i\alpha}b_\alpha V_{j\alpha}^*.
\end{align*}
特别地 $C_{ii}=0$。因此三阶行列式只剩两个三循环项：
\begin{equation*}
 \det C=C_{12}C_{23}C_{31}+C_{13}C_{21}C_{32}.
\end{equation*}
定义
\begin{equation*}
 \Delta_u\equiv
 (a_1-a_2)(a_1-a_3)(a_2-a_3).
\end{equation*}
第一个循环的上型谱因子为
\begin{align*}
 &(a_1-a_2)(a_2-a_3)(a_3-a_1)\\
 &=-\Delta_u,
\end{align*}
第二个循环的上型谱因子为
\begin{align*}
 &(a_1-a_3)(a_2-a_1)(a_3-a_2)\\
 &=+\Delta_u.
\end{align*}
又因为 $H_d$ 厄米，令
\begin{equation*}
 z\equiv(H_d)_{12}(H_d)_{23}(H_d)_{31},
\end{equation*}
则
\begin{equation*}
 (H_d)_{13}(H_d)_{21}(H_d)_{32}=z^*.
\end{equation*}
所以
\begin{align*}
 \det C
 &=\Delta_u(z^*-z)\\
 &=-2\ii\Delta_u\Im z.
\end{align*}
问题于是化为计算 $\Im z$。

由于单位矩阵的非对角元为零，给 $H_d$ 加上任意常数倍单位矩阵都不改变 $z$。因此可以从所有 $b_\alpha$ 中减去 $b_3$，定义
\begin{equation*}
 x\equiv b_1-b_3,
 \qquad
 y\equiv b_2-b_3.
\end{equation*}
对 $i\neq j$，
\begin{equation*}
 (H_d)_{ij}
 =xV_{i1}V_{j1}^*+yV_{i2}V_{j2}^*.
\end{equation*}
于是 $z$ 是 $x,y$ 的三次齐次多项式。纯 $x^3$ 项为
\begin{equation*}
 x^3|V_{11}V_{21}V_{31}|^2,
\end{equation*}
纯 $y^3$ 项为
\begin{equation*}
 y^3|V_{12}V_{22}V_{32}|^2,
\end{equation*}
二者都为实数，不贡献 $\Im z$。

$x^2y$ 的三个交叉项的虚部分别是
\begin{align*}
 \Im\!\big[(V_{11}V_{21}^*)(V_{21}V_{31}^*)(V_{32}V_{12}^*)\big]
 &=-J_{\rm CKM}|V_{21}|^2,\\
 \Im\!\big[(V_{11}V_{21}^*)(V_{22}V_{32}^*)(V_{31}V_{11}^*)\big]
 &=-J_{\rm CKM}|V_{11}|^2,\\
 \Im\!\big[(V_{12}V_{22}^*)(V_{21}V_{31}^*)(V_{31}V_{11}^*)\big]
 &=-J_{\rm CKM}|V_{31}|^2.
\end{align*}
三代幺正矩阵的四元积满足
\begin{equation*}
 \Im(V_{i\alpha}V_{j\beta}V_{i\beta}^*V_{j\alpha}^*)
 =J_{\rm CKM}\,\epsilon_{ijk}\epsilon_{\alpha\beta\gamma},
\end{equation*}
其中未写出的 $k,\gamma$ 是各自剩余的第三个指标。由第一列归一化
\begin{equation*}
 |V_{11}|^2+|V_{21}|^2+|V_{31}|^2=1,
\end{equation*}
故全部 $x^2y$ 项的虚部为
\begin{equation*}
 -J_{\rm CKM}x^2y.
\end{equation*}
对 $xy^2$ 项逐个应用四元积恒等式，三个虚部分别是
\begin{align*}
 \Im(V_{11}V_{22}V_{12}^*V_{21}^*)|V_{32}|^2
 &=+J_{\rm CKM}|V_{32}|^2,\\
 \Im(V_{21}V_{32}V_{22}^*V_{31}^*)|V_{12}|^2
 &=+J_{\rm CKM}|V_{12}|^2,\\
 \Im(V_{31}V_{12}V_{32}^*V_{11}^*)|V_{22}|^2
 &=+J_{\rm CKM}|V_{22}|^2.
\end{align*}
三项相加后利用第二列归一化得到
\begin{equation*}
 +J_{\rm CKM}xy^2.
\end{equation*}
因此
\begin{align*}
 \Im z
 &=J_{\rm CKM}xy(y-x)\\
 &=-J_{\rm CKM}(b_1-b_2)(b_1-b_3)(b_2-b_3)\\
 &=-J_{\rm CKM}\Delta_d.
\end{align*}
代回 $\det C=-2\ii\Delta_u\Im z$：
\begin{equation*}
 \det[H_u,H_d]
 =2\ii J_{\rm CKM}\Delta_u\Delta_d,
\end{equation*}
即\eqnrefs{eq:jarlskog-determinant-dp11}。

这个推导还直接给出两个退化检查：若任意两种上型 Yukawa 本征值相同，则 $\Delta_u=0$；若任意两种下型本征值相同，则 $\Delta_d=0$。任一情形都使所有弱基不变的 $CP$ 奇量消失，因为相应简并子空间内出现额外的味旋转自由度。
\end{derivation}
